
\begin{table}[htbp]
   \caption{\label{tab:education} Education Heterogeneity: DDD Estimates by Education Level}
   \bigskip
   \centering
   \begin{adjustbox}{max width=\textwidth}
   \begin{tabular}{lcccc}
      \toprule
       & \multicolumn{2}{c}{New Hire Rate} & \multicolumn{2}{c}{Separation Rate}\\
                                                            & (1)           & (2)           & (3)           & (4)\\  
      \midrule 
      expansion\_state $\times$ post $\times$ high\_esi     & 0.5296$^{**}$ & 0.8524$^{**}$ & 0.4373$^{*}$  & 0.8981$^{**}$\\   
                                                            & (0.2153)      & (0.3686)      & (0.2546)      & (0.4094)\\   
       \\
      Observations                                          & 4,122         & 4,122         & 4,122         & 4,122\\  
      R$^2$                                                 & 0.99021       & 0.98964       & 0.99088       & 0.99274\\  
       \\
      state\_ind fixed effects                              & $\checkmark$  & $\checkmark$  & $\checkmark$  & $\checkmark$\\   
      ind\_time fixed effects                               & $\checkmark$  & $\checkmark$  & $\checkmark$  & $\checkmark$\\   
      state\_time fixed effects                             & $\checkmark$  & $\checkmark$  & $\checkmark$  & $\checkmark$\\   
      \bottomrule
   \end{tabular}
   \end{adjustbox}
   
   \par \raggedright 
   Columns (1)--(2) report DDD estimates for new hire rate; columns (3)--(4) for separation rate.\\
   Low education: less than bachelor's degree. High education: bachelor's or above.\\
   Medicaid expansion primarily affected low-education workers; high-education serves as within-industry placebo.\\
   All specifications include state$\times$industry, industry$\times$time, and state$\times$time FE.\\
   Standard errors clustered at state level. * p$<$0.10, ** p$<$0.05, *** p$<$0.01.
\end{table}


